\documentclass[11pt,a4paper]{article}
\usepackage[utf8]{inputenc}
\usepackage[german]{babel}
\usepackage[margin=2.5cm]{geometry}
\usepackage{graphicx}
\usepackage{xcolor}
\usepackage{fancyhdr}
\usepackage{titlesec}
\usepackage{hyperref}
\usepackage{array}

% Farben definieren
\definecolor{primary}{RGB}{0, 102, 204}
\definecolor{secondary}{RGB}{51, 51, 51}
\definecolor{accent}{RGB}{255, 153, 0}

% Seitenstil
\pagestyle{fancy}
\fancyhf{}
\fancyfoot[C]{\thepage}
\renewcommand{\headrulewidth}{0pt}

% Titel-Formatierung
\titleformat{\section}{\Large\bfseries\color{primary}}{}{0em}{}[\color{primary}\rule{\linewidth}{2pt}]
\titleformat{\subsection}{\large\bfseries\color{secondary}}{}{0em}{}

\title{}
\author{}
\date{}

\begin{document}

% ===== TITELSEITE =====
\thispagestyle{empty}
\vspace*{3cm}

\begin{center}
    {\fontsize{48}{60}\selectfont\bfseries\color{primary}BUSINESSPLAN}
    
    \vspace{0.5cm}
    {\fontsize{24}{30}\selectfont\color{accent}[Ihr Startup-Name]}
    
    \vspace{3cm}
    {\fontsize{14}{18}\selectfont\color{secondary}
    Ein umfassendes Geschäftskonzept für die Gründung und Entwicklung\\
    unseres innovativen Unternehmens
    }
\end{center}

\vspace{4cm}

\begin{center}
    \begin{tabular}{l l}
        \textbf{Gründer:} & [Name des Gründers] \\[0.3cm]
        \textbf{Adresse:} & [Adresse] \\[0.3cm]
        \textbf{Telefon:} & [Telefonnummer] \\[0.3cm]
        \textbf{E-Mail:} & [E-Mail-Adresse] \\[0.3cm]
        \textbf{Datum:} & \today
    \end{tabular}
\end{center}

\vfill

\begin{center}
    {\color{primary}\rule{0.8\linewidth}{3pt}}
    
    \vspace{0.5cm}
    {\fontsize{11}{13}\selectfont\color{secondary}
    Dieses Dokument enthält vertrauliche Informationen.\\
    Unbefugte Weitergabe ist untersagt.
    }
\end{center}

\newpage

% ===== INHALTSVERZEICHNIS =====
\tableofcontents
\newpage

% ===== 1. EXECUTIVE SUMMARY =====
\section{Executive Summary}
Prägnante Zusammenfassung des Geschäftskonzepts (max. 1-2 Seiten).
\begin{itemize}
    \item Geschäftsidee kurz beschreiben
    \item Zielmarkt und Marktpotenzial
    \item Wettbewerbsvorteil
    \item Finanzielle Prognose (Umsatz, Gewinn)
    \item Kapitalbedarfe
\end{itemize}

% ===== 2. UNTERNEHMENSKONZEPT =====
\section{Unternehmenskonzept}

\subsection{Geschäftsidee}
Detaillierte Beschreibung der Geschäftsidee und des Produkts/der Dienstleistung.

\subsection{Vision und Mission}
\begin{itemize}
    \item \textbf{Vision:} Langfristige Zielstellung
    \item \textbf{Mission:} Unternehmenszweck und Werte
\end{itemize}

\subsection{Rechtsform und Organisationsstruktur}
Wahl der Unternehmensform (GmbH, GbR, etc.) und interne Struktur.

% ===== 3. MARKTANALYSE =====
\section{Marktanalyse}

\subsection{Marktgröße und Markttrends}
\begin{itemize}
    \item Gesamtmarktvolumen
    \item Wachstumstrends
    \item Prognosen für die nächsten 5 Jahre
\end{itemize}

\subsection{Zielgruppe}
\begin{itemize}
    \item Demografie
    \item Psychografie
    \item Kundenbedarf
\end{itemize}

\subsection{Wettbewerbsanalyse}
\begin{itemize}
    \item Direkte Konkurrenten
    \item Indirekte Konkurrenten
    \item Wettbewerbsvorteil (USP)
    \item SWOT-Analyse
\end{itemize}

% ===== 4. MARKETING- UND VERTRIEBSSTRATEGIE =====
\section{Marketing- und Vertriebsstrategie}

\subsection{Produktpositionierung}
Wie soll das Produkt/die Dienstleistung positioniert werden?

\subsection{Marketingmaßnahmen}
\begin{itemize}
    \item Online-Marketing (SEO, Social Media, etc.)
    \item Offline-Marketing (Events, PR, etc.)
    \item Partnerschaften und Kooperationen
\end{itemize}

\subsection{Vertriebskanäle}
Beschreibung der Vertriebswege (direkt, indirekt, online, etc.).

\subsection{Preisgestaltung}
Pricing-Strategie und Begründung.

% ===== 5. LEISTUNGSERSTELLUNG =====
\section{Leistungserstellung}

\subsection{Produktentwicklung}
Entwicklungsprozess und Roadmap.

\subsection{Lieferanten und Partnerschaften}
\begin{itemize}
    \item Kritische Lieferanten
    \item Qualitätssicherung
\end{itemize}

\subsection{Technologie und Infrastruktur}
Benötigte Systeme, Software, Hardware.

% ===== 6. MANAGEMENT UND ORGANISATION =====
\section{Management und Organisation}

\subsection{Organisationsstruktur}
Aufbauorganisation und Verantwortlichkeiten.

\subsection{Gründerteam}
\begin{itemize}
    \item Kurzbios der Gründer
    \item Relevant Erfahrung und Expertise
    \item Referenzen
\end{itemize}

\subsection{Personalplanung}
\begin{itemize}
    \item Aktueller und zukünftiger Personalbestand
    \item Qualifikationsanforderungen
\end{itemize}

% ===== 7. FINANZPLANUNG =====
\section{Finanzplanung}

\subsection{Kapitalbedarfsplan}
Detaillierung des benötigten Kapitals.

\subsection{Rentabilitätsprognose}
\begin{itemize}
    \item Umsatzprognose (3-5 Jahre)
    \item Kostenstruktur
    \item Gewinnprognose
\end{itemize}

\subsection{Break-Even-Analyse}
Zeitpunkt der Profitabilität.

\subsection{Finanzierungsquellen}
\begin{itemize}
    \item Eigenkapital
    \item Fremdkapital (Kredite, Investoren)
\end{itemize}

% ===== 8. RISIKOMANAGEMENT =====
\section{Risikomanagement}

\subsection{Identifizierte Risiken}
\begin{itemize}
    \item Marktrisiken
    \item Technologische Risiken
    \item Personalrisiken
    \item Finanzielle Risiken
\end{itemize}

\subsection{Risikominderungsmaßnahmen}
Strategien zur Risikovermeidung und -minimierung.

% ===== 9. MEILENSTEINE UND ZEITPLAN =====
\section{Meilensteine und Zeitplan}

Konkrete Ziele und Zeitrahmen für die ersten 24 Monate.

% ===== 10. ANHANG =====
\section{Anhang}

\begin{itemize}
    \item Detaillierte Finanzpläne (Bilanz, GuV, Cashflow)
    \item Marktforschungsergebnisse
    \item Produktmuster/Prototypen
    \item Lebenslauf der Gründer
    \item Zertifikate und Qualifikationen
    \item Presseberichte und Referenzen
\end{itemize}

\end{document}